\documentclass[10pt,a4paper]{amsart}
\usepackage{amsmath,amssymb,amsthm,graphicx}
\usepackage{verbatim}
\newenvironment{reference}{\comment}{\endcomment}
%\newenvironment{reference}{}{}
\newenvironment{slogan}{\comment}{\endcomment}
\newenvironment{history}{\comment}{\endcomment}

% For commutative diagrams we use Xy-pic
\usepackage[all]{xy}

% We use 2cell for 2-commutative diagrams.
\xyoption{2cell}
% Theorem environments.
%
\theoremstyle{plain}
\newtheorem{theorem}[subsection]{Theorem}
\newtheorem{proposition}[subsection]{Proposition}
\newtheorem{lemma}[subsection]{Lemma}

\theoremstyle{definition}
\newtheorem{definition}[subsection]{Definition}
\newtheorem{example}[subsection]{Example}
\newtheorem{exercise}[subsection]{Exercise}
\newtheorem{situation}[subsection]{Situation}

\theoremstyle{remark}
\newtheorem{remark}[subsection]{Remark}
\newtheorem{remarks}[subsection]{Remarks}

\numberwithin{equation}{subsection}

% Macros
%
\def\lim{\mathop{\mathrm{lim}}\nolimits}
\def\colim{\mathop{\mathrm{colim}}\nolimits}
\def\Spec{\mathop{\mathrm{Spec}}}
\def\Hom{\mathop{\mathrm{Hom}}\nolimits}
\def\Ext{\mathop{\mathrm{Ext}}\nolimits}
\def\SheafHom{\mathop{\mathcal{H}\!\mathit{om}}\nolimits}
\def\SheafExt{\mathop{\mathcal{E}\!\mathit{xt}}\nolimits}
\def\Sch{\mathit{Sch}}
\def\Mor{\mathop{\mathrm{Mor}}\nolimits}
\def\Ob{\mathop{\mathrm{Ob}}\nolimits}
\def\Sh{\mathop{\mathit{Sh}}\nolimits}
\def\NL{\mathop{N\!L}\nolimits}
\def\CH{\mathop{\mathrm{CH}}\nolimits}
\def\proetale{{pro\text{-}\acute{e}tale}}
\def\etale{{\acute{e}tale}}
\def\QCoh{\mathit{QCoh}}
\def\Ker{\mathop{\mathrm{Ker}}}
\def\Im{\mathop{\mathrm{Im}}}
\def\Coker{\mathop{\mathrm{Coker}}}
\def\Coim{\mathop{\mathrm{Coim}}}

% Boxtimes
%
\DeclareMathSymbol{\boxtimes}{\mathbin}{AMSa}{"02}

%
% Macros for moduli stacks/spaces
%
\def\QCohstack{\mathcal{QC}\!\mathit{oh}}
\def\Cohstack{\mathcal{C}\!\mathit{oh}}
\def\Spacesstack{\mathcal{S}\!\mathit{paces}}
\def\Quotfunctor{\mathrm{Quot}}
\def\Hilbfunctor{\mathrm{Hilb}}
\def\Curvesstack{\mathcal{C}\!\mathit{urves}}
\def\Polarizedstack{\mathcal{P}\!\mathit{olarized}}
\def\Complexesstack{\mathcal{C}\!\mathit{omplexes}}
% \Pic is the operator that assigns to X its picard group, usage \Pic(X)
% \Picardstack_{X/B} denotes the Picard stack of X over B
% \Picardfunctor_{X/B} denotes the Picard functor of X over B
\def\Pic{\mathop{\mathrm{Pic}}\nolimits}
\def\Picardstack{\mathcal{P}\!\mathit{ic}}
\def\Picardfunctor{\mathrm{Pic}}
\def\Deformationcategory{\mathcal{D}\!\mathit{ef}}
\usepackage{cleveref}
\crefname{originalreference}{original reference}{original references}
\Crefname{originalreference}{Original reference}{Original references}
\makeatletter
\newlabel{lemma-algebraically-closed}{{09GR}{source}{}{sourcepointer-1}{}}
\newlabel{lemma-algebraically-closed@cref}{{[originalreference][0][]09GR}{[1][0][]source}{}{}{}}
\makeatother
\crefname{originalreference}{original reference}{original references}
\Crefname{originalreference}{Original reference}{Original references}
\makeatletter
\newlabel{lemma-map-into-algebraic-closure}{{09GU}{source}{}{sourcepointer-2}{}}
\newlabel{lemma-map-into-algebraic-closure@cref}{{[originalreference][0][]09GU}{[1][0][]source}{}{}{}}
\makeatother
\title{Normality is the constancy of embedding images}
\author{The Stacks Project Authors}
\date{stacks-project:fields.tex:local-snapshot-20261012; GNU FDL 1.2 or later}
\begin{document}
\maketitle
\noindent\textit{Editorial note.} All delivered mathematical text is the author's own, copied from the pinned Stacks Project chapter file fields.tex; no statement or proof step was rewritten, reordered or removed.\par\medskip
\noindent\textit{Editorial note.} The counted proof is the author's complete argument: This characterization is the workhorse behind the lift-maps lemma (delivered L48358) and the automorphism-counting lemma (delivered L48349).\par\medskip
\noindent\textit{Editorial note.} The label lemma-characterize-normal names the counted statement, and the proof environment immediately following it in the source is delivered verbatim as its proof.\par\medskip
\noindent\textit{Editorial note.} The counted statement is carried by the author's own lemma environment in the chapter source.\par\medskip
\noindent\textit{Allowed standard prerequisites (editorial).}\begin{itemize}
\item Chapter definition of normal: an algebraic extension E over F is normal when every irreducible polynomial over F with a root in E splits completely over E
\item Chapter lemma-map-into-algebraic-closure (delivered card L48339): F-maps of algebraic extensions extend into a chosen algebraic closure
\item Chapter lemma-algebraically-closed (delivered card L48337): over an algebraically closed field every polynomial splits into linear factors
\item Isomorphism invariance: if two subfields of the algebraic closure are F-isomorphic, an element and its image have the same minimal polynomial over F
\end{itemize}\medskip
\ifdefined\hypertarget\hypertarget{sourcepointer-1}{}\fi
\noindent\textit{Source reference.} Original reference 09GR: sources/stacks-fields/pinned-original-source.tex, source lines 1128--1128. Prerequisite pointer invoked at 2132; it names a standard fact of the same chapter that the counted argument uses, and every step of the counted proof built on it is retained verbatim in the delivered body.\par\medskip
\ifdefined\hypertarget\hypertarget{sourcepointer-2}{}\fi
\noindent\textit{Source reference.} Original reference 09GU: sources/stacks-fields/pinned-original-source.tex, source lines 1240--1240. Prerequisite pointer invoked at 2122; it names a standard fact of the same chapter that the counted argument uses, and every step of the counted proof built on it is retained verbatim in the delivered body.\par\medskip
\begin{lemma}
\label{lemma-characterize-normal}
Let $E/F$ be an algebraic extension of fields. Let $\overline{F}$ be an
algebraic closure of $F$. The following are equivalent
\begin{enumerate}
\item $E$ is normal over $F$, and
\item for every pair $\sigma, \sigma' \in \Mor_F(E, \overline{F})$ we
have $\sigma(E) = \sigma'(E)$.
\end{enumerate}
\end{lemma}

\begin{proof}
Let $\mathcal{P}$ be the set of all minimal polynomials over $F$ of
all elements of $E$. Set
$$
T =
\{\beta \in \overline{F} \mid P(\beta) = 0\text{ for some }P \in \mathcal{P}\}
$$
It is clear that if $E$ is normal over $F$, then $\sigma(E) = T$
for all $\sigma \in \Mor_F(E, \overline{F})$. Thus we see that (1)
implies (2).

\medskip\noindent
Conversely, assume (2). Pick $\beta \in T$.
We can find a corresponding $\alpha \in E$ whose minimal polynomial
$P \in \mathcal{P}$ annihilates $\beta$. Because $F(\alpha) = F[x]/(P)$
we can find an element $\sigma_0 \in \Mor_F(F(\alpha), \overline{F})$ mapping
$\alpha$ to $\beta$. By Lemma \ref{lemma-map-into-algebraic-closure}
we can extend $\sigma_0$ to a $\sigma \in \Mor_F(E, \overline{F})$.
Whence we see that $\beta$ is in the common image of all embeddings
$\sigma : E \to \overline{F}$. It follows that $\sigma(E) = T$
for any $\sigma$. Fix a $\sigma$. Now let $P \in \mathcal{P}$. Then we
can write
$$
P = (x - \beta_1) \ldots (x - \beta_n)
$$
for some $n$ and $\beta_i \in \overline{F}$ by
Lemma \ref{lemma-algebraically-closed}. Observe that $\beta_i \in T$.
Thus $\beta_i = \sigma(\alpha_i)$ for some $\alpha_i \in E$. Thus
$P = (x - \alpha_1) \ldots (x - \alpha_n)$ splits completely over $E$.
This finishes the proof.
\end{proof}

\end{document}
